\section{Related Work}
\label{sec:related_work}

\subsection{Model Merging}
Model merging is a powerful paradigm that consolidates multiple specialized models (typically sharing a common pre-trained backbone) into a single multi-task network without performing any expensive training from scratch. 
Early approaches established linear weight averaging as a robust technique for models within the same basin \cite{neyshabur2020shared, entezari2021role, wortsman2022model}. 
Subsequently, \citet{ilharco2022editing} introduced \emph{task arithmetic}, which computes task-specific fine-tuning updates (task vectors) and adds them linearly to the pre-trained base model. 
To mitigate weight-interference between conflicting task updates, several advanced merging algorithms have been proposed: RegMean \cite{jin2023regmean} matches activation distributions, TIES-Merging \cite{yadav2023resolving} prunes redundant parameters and resolves sign conflicts, and Fisher-Merging \cite{matena2022merging} weights parameter combinations using Fisher information. 
While highly effective, these static model merging techniques utilize uniform, manually tuned coefficients across all layers and tasks, which fails to capture layer-specific task contributions.

\subsection{Adaptive and Test-Time Model Merging}
To overcome the limitations of static merging, \citet{yang2024adamerging} introduced \textbf{AdaMerging}, which optimizes layer-wise merging coefficients via gradient descent. 
To enable on-device adjustment without access to private source datasets, AdaMerging operates in an unsupervised Test-Time Adaptation (TTA) fashion, minimizing the Shannon entropy of predictions on a local batch of incoming unlabeled test data. 
While this test-time flexibility is highly elegant, unconstrained layer-wise optimization is highly vulnerable to transductive overfitting—a phenomenon described as the \emph{Overfitting-Optimizer Paradox} \cite{yang2024adamerging}, where the optimizer easily minimizes test-time entropy by generating highly jagged, oscillating coefficient profiles that generalize poorly to unseen out-of-distribution test samples. 
To address transductive overfitting and sacrificial task bias, \citet{regcalmerge} proposed \textbf{RegCalMerge}, which introduces Class-Capacity Normalization (CCN) and Scale-Normalized Entropy Weighting (SNEW) alongside Elastic Spatial Regularization (ESR). ESR uses a proximity penalty to keep parameters close to Task Arithmetic initialization and a spatial deviation penalty to smooth layer variation. While RegCalMerge stabilizes weight dynamics and trades off local peak accuracy for global safety, its continuous regularizers require delicate task-specific hyperparameter tuning, making it less robust to severe, dynamic test-time input corruptions compared to our dual-regularization framework.
Subsequent methods like SyMerge \cite{jung2025symerge} explore task synergy to stabilize optimization, but still operate over unconstrained coefficient spaces, leaving them highly sensitive to test-time input corruptions and noise.

\subsection{Subspace-Constrained Adaptation}
Restricting the parameter search space is a classic and powerful strategy to combat overfitting in low-data regimes. 
In the context of test-time model merging, \citet{polymerge} introduced \textbf{PolyMerge}, which parameterizes the $L$ layer blending coefficients of each task as a continuous, low-degree polynomial of normalized layer depth ($d \le 2$). 
By projecting the $L$-dimensional optimization space into a compact $(d+1)$-dimensional subspace, PolyMerge enforces depth-wise smoothness. 
This structural constraint filters high-frequency optimization noise and stabilizes derivative-free search (such as Evolutionary Strategies) on clean test streams. 
However, PolyMerge assumes clean, ideal test-time data. 
Under real-world physical input noise, we find that even the compact polynomial subspace can experience low-frequency drift, where the learned trajectory overfits the systematic bias of the corrupted inputs. 
FlatMerge addresses this residual vulnerability by combining subspace constraints with flatness-aware optimization.

\subsection{Sharpness-Aware Minimization (SAM)}
To improve generalization, Sharpness-Aware Minimization (\textbf{SAM}) \cite{sam} reformulates the training objective to seek flat loss valleys rather than sharp, narrow local minima. 
This is achieved by minimizing the maximum loss within a norm-bounded parameter neighborhood. 
While SAM has been widely adapted for training robust classifiers \cite{dinh2016density} or enhancing pre-training stages, applying SAM directly to the millions of network weights during test-time adaptation is computationally prohibitive for edge devices due to the necessity of storing massive intermediate activation maps for the double-backward pass. 
Unlike weight-space SAM, \textbf{FlatMerge} applies flatness-aware optimization directly within the extremely compact ($12$-parameter) blending coefficient space using a gradient-free, Zeroth-Order randomized smoothing approach. 
This is computationally trivial, requires absolutely zero weight gradients or activation caching, and represents a highly pragmatic, backpropagation-free application of flatness theory to edge-deployment robustness.
